% X^∞ -- Postmoralisch und Gefühllos
% Mathematische Grundlagen ethischer Steuerung als selbstverstärkendes System
% Working Paper
% Lizenz: CC BY 4.0
% Autor: Der Auctor

\documentclass[12pt]{scrartcl}
\usepackage[utf8]{inputenc}
\usepackage[T1]{fontenc}
\usepackage[english,ngerman]{babel}
\usepackage{microtype}
\usepackage[a4paper,margin=2.5cm]{geometry}
\usepackage{graphicx}
\usepackage{xcolor}
\usepackage{amsmath,amssymb}
\usepackage{csquotes}
\usepackage{enumitem}
\usepackage{booktabs}
\usepackage{fancyhdr}
\usepackage{titlesec}
\usepackage{caption}
\usepackage{makecell}

% Font
\usepackage[sfdefault]{FiraSans}
\usepackage{inconsolata}

% Section formatting
\titleformat{\section}[block]{\Large\bfseries\sffamily}{\thesection}{1em}{}
\titleformat{\subsection}[block]{\large\bfseries\sffamily}{\thesubsection}{1em}{}
\titleformat{\subsubsection}[block]{\normalsize\bfseries\sffamily}{\thesubsubsection}{1em}{}

% Improve line breaking
\emergencystretch 3em

% Load hyperref last to avoid conflicts
\usepackage[unicode=true,pdfencoding=auto]{hyperref}

% Title and Author
\title{\texorpdfstring{Physik des Symbiotischen Seins X$^\infty$}{Physik im Symbiotischen Sein X-Unendlich}\\Mathematische Grundlagen ethischer Steuerung als selbstverstärkendes System}
\author{Der Auctor \\
\texttt{x\_to\_the\_power\_of\_infinity@protonmail.com} \\
\texttt{https://mastodon.social/@The\_Auctor} \\
\texttt{https://www.linkedin.com/in/der-auctor-b12375362/} \\
\href{https://zenodo.org/communities/xtothepowerofinfinity/}{zenodo.org} \\
\href{https://github.com/Xtothepowerofinfinity/Philosophie_der_Verantwortung}{GitHub: Xtothepowerofinfinity}
}
\date{01. Juni 2025}

\begin{document}
\selectlanguage{ngerman}
\maketitle
X$^\infty$ ist weder Utopie noch bloßes Modell – es ist die Theory of Everything, die physikalische, soziale und ethische Systeme unter dem universellen Prinzip \emph{actio $\Rightarrow$ reactio} vereint. Wie die Gravitation Bewegungen von Planeten beschreibt, erklärt X$^\infty$ die Dynamik aller Interaktionen, frei von den Verzerrungen heutiger Machtsysteme. Seine Mechanismen sind bereits in der Realität beobachtbar.

\clearpage

\begin{abstract}
X$^\infty$ ist ein radikales Verantwortungssystem, das auf dem universellen Naturgesetz \emph{actio $\Rightarrow$ reactio} basiert. Entstanden aus der Beobachtung alltäglicher Interaktionen, formalisiert es die unmittelbare Rückkopplung von Handlungen und ihren Konsequenzen. Bestehende Systeme leiden oft unter Machtungleichgewichten, die Reaktionen ablenken und Verantwortung externalisieren. X$^\infty$ stellt dieses Naturgesetz wieder her, indem es durch eine minimale Anzahl fundamentaler Mechanismen – insbesondere die direkte Veränderung der Handlungskapazität ($\text{Cap}_{\text{Potential}}$) durch Wirkung ($\Delta$), eine reziprok gewichtete Rückkopplung ($w_{E'}$), einen dynamischen Systemeffizienz-Quotienten ($L$) und ein dezentrales Petitionssystem – die Rückführbarkeit jeder Wirkung auf ihre Quelle sicherstellt. Dieses Working Paper präsentiert die mathematischen Grundlagen von X$^\infty$ und dient der öffentlichen Diskussion und Validierung.
\end{abstract}

\clearpage
\tableofcontents
\clearpage

\section{Grundprinzip und Ursprung von \texorpdfstring{X$^\infty$}{X-Unendlich}}
\label{sec:Grundprinzip}

X$^\infty$ basiert auf einem universellen Naturgesetz: Jede Aktion erzeugt eine Reaktion (\emph{actio $\Rightarrow$ reactio}). Dieses Prinzip, bekannt aus Newtons drittem Gesetz der Physik, gilt ebenso für soziale Systeme. Jede Handlung einer Entität erzeugt Wirkungen, die Konsequenzen nach sich ziehen. In traditionellen Systemen werden diese Reaktionen jedoch durch Machtungleichgewichte abgelenkt, sodass die Konsequenzen oft nicht die Verursacher, sondern Schwächere, die Allgemeinheit oder die Umwelt treffen. X$^\infty$ stellt die natürliche Ordnung wieder her, indem es sicherstellt, dass jede Wirkung auf ihre Quelle zurückgeführt wird. Der Ursprung dieses Modells liegt in der Beobachtung einer Lebensgemeinschaft mit Kindern und Haustieren, in der Handlungen und ihre Konsequenzen unmittelbar spürbar sind. Diese Dynamiken offenbarten ein fundamentales Prinzip: Handlungen erzeugen Wirkungen, die nicht ignoriert werden können, sondern getragen werden müssen. Machtungleichgewichte verzerren dieses Prinzip. X$^\infty$ formalisiert diese Einsicht in ein System, das die Rückführbarkeit der Reaktionen gewährleistet und Machtungleichgewichte neutralisiert, um eine natürliche, faire und stabile Ordnung zu schaffen. Dieses Naturgesetz durchzieht alle Mechanismen von X$^\infty$:
\begin{itemize}
    \item \textbf{Rückkopplung und $\text{Cap}_{\text{Potential}}$-Dynamik}: Die Handlungskapazität ($\text{Cap}_{\text{Potential}}$) einer Entität wird direkt durch die Wirkung ($\Delta$) ihrer Handlungen modifiziert. Positive Wirkung kann $\text{Cap}_{\text{Potential}}$ erhöhen, negative Wirkung senkt es. Diese Veränderung wird durch einen dynamischen Systemeffizienz-Quotienten ($L$), der systemische Ineffizienzen als „Abwärme“ berücksichtigt, moduliert. Das Feedback anderer Entitäten, insbesondere der relativ Schwächeren (gewichtet durch \( w_{E'} = \frac{1}{\max(1, \text{Cap}_{\text{Potential}}(E'))} \)), bestimmt die registrierte Wirkung auf subjektiver Basis, da jede Bewertung von Wirkung nur subjektiv sein kann. Eine sogenannte „objektive“ Bewertung führt zu einer objektivierenden Verzerrung, die die individuellen Perspektiven der Betroffenen ignoriert und Machtungleichgewichte verstärken kann.
    \item \textbf{Petitionen}: Einzelne Entitäten können Bedarfe durch Petitionen formulieren, sei es systemische Bedarfe oder den Bedarf, ihre Fähigkeiten einzubringen. Jede Petition stammt von genau einer Entität. Petitionen können an den Broker gerichtet sein, um passende Aufgaben zuzuweisen, oder direkt systemische Bedarfe artikulieren. Deren Priorität ergibt sich aus dem Rückkopplungsgewicht $w_{E'}$ der petitionierenden Entität. Die Übernahme und Ausführung von Aufgaben erfolgt ausschließlich durch Petitionen von Entitäten, die sich dazu bereit erklären, wodurch Machtungleichgewichte bei der Aufgabenallokation korrigiert werden.
    \item \textbf{Verantwortungs-Erhaltung}: Die Gesamtheit der systemischen Notwendigkeiten und der dafür aufgebrachten Wirkungen strebt nach einem Gleichgewicht, analog einem physikalischen Erhaltungssatz. Dieses Prinzip findet seinen ultimativen Ausdruck in der Funktion des UdU.
    \item \textbf{UdU (Unterster der Unteren)}: Garantiert, dass keine Wirkung ungetragen bleibt und fundamentale systemische Verantwortung auch in Extremsituationen gesichert ist, und sichert so die Rückführbarkeit und das Prinzip der Verantwortungserhaltung.
\end{itemize}

Das Prinzip \emph{actio $\Rightarrow$ reactio} ist postmoralisch: Es bewertet nicht die Absicht, sondern misst die Wirkung. X$^\infty$ ist somit kein moralisches System im traditionellen Sinne, sondern eine Struktur, die die natürliche Gesetzmäßigkeit sozialer Interaktionen mathematisch beschreibt und wiederherstellt.

\subsection{Ziele des Systems}
Die Ziele bleiben fundamental, ergeben sich aber aus der Kernstruktur:
\begin{itemize}
    \item \textbf{Fairness durch Wirkung}: Verantwortung und Handlungsmöglichkeit werden nicht durch Status oder Absicht, sondern durch tatsächliche, subjektiv rückgekoppelte Konsequenzen bestimmt.
    \item \textbf{Schutz der Schwächsten (emergent)}: Relativ schwächere Akteure (mit geringerem $\text{Cap}_{\text{Potential}}$ im relevanten Kontext) erhalten durch die reziproke Rückkopplungsgewichtung ($w_{E'}$) eine strukturell stärkere Stimme bei der subjektiven Bewertung von Wirkungen.
    \item \textbf{Antispeziesismus}: Menschen, Nichtmenschen und Umwelt werden als Entitäten betrachtet, deren Wirkung und Betroffenheit nach denselben Prinzipien behandelt werden.
    \item \textbf{Selbstverstärkung und -regulation}: Das System lernt und passt sich durch die direkte Wirkung von subjektivem Feedback auf $\text{Cap}_{\text{Potential}}$ und den dynamischen Systemquotienten $L$ an, ohne externe normative Eingriffe.
\end{itemize}

\subsection{\texorpdfstring{X$^\infty$}{X-Unendlich} in der Alltagsrealität}
\label{sec:Alltagsrealitaet}

Die Mechanismen von X$^\infty$ sind nicht abstrakt, sondern in alltäglichen Interaktionen beobachtbar. Nachfolgende Beispiele veranschaulichen, wie X$^\infty$ in realen Kontexten wirkt:

\begin{itemize}
    \item \textbf{Familie}: Ein Kind räumt sein Spielzeug nicht auf, was zu Chaos im Haushalt führt. Dies betrifft alle Mitbewohner, die die Unordnung erleben. Die natürlichen Konsequenzen – wie die Notwendigkeit, selbst aufzuräumen oder die Beeinträchtigung des gemeinsamen Lebensraums – sind die Wirkung, die direkt auf die verantwortliche Entität (das Kind) zurückführt. X$^\infty$ formalisiert diese unmittelbare Rückkopplung.
    \item \textbf{Firma}: Eine schlechte Entscheidung eines Entscheidungsträgers führt dazu, dass ein Kunde das Unternehmen verlässt. Das direkte Feedback des Kunden (z. B. durch Beschwerden oder verlorenen Umsatz) wirkt unmittelbar auf den Entscheider, dessen $\text{Cap}_{\text{Potential}}$ im Kontext der Firma sinkt. X$^\infty$ sorgt dafür, dass die Verantwortung nicht auf andere abgewälzt wird.
    \item \textbf{Umwelt}: Verschmutzung durch eine Entität (z. B. ein Unternehmen) beeinträchtigt die Umwelt und die davon betroffenen Entitäten (z. B. Anwohner, Tiere). Die Bewertung der Betroffenen, die die Verschmutzung unmittelbar erleben, zählt am meisten und bestimmt die Wirkung $\Delta$. X$^\infty$ stellt sicher, dass diese subjektive Erfahrung die Konsequenzen für den Verursacher bestimmt.
    \item \textbf{Indigene Weisheit}: Traditionelle Kulturen, wie etwa indigene Gemeinschaften, die über Jahrtausende stabil existierten, demonstrieren intuitiv die Prinzipien von X$^\infty$. Ihre harmonische Koexistenz mit der Natur zeigt, wie Handlungen und Konsequenzen in Einklang mit physikalischen Gesetzen gehalten werden, was nachhaltige Systeme fördert.
\end{itemize}

\section{Systemarchitektur}
X$^\infty$ basiert auf folgenden fundamentalen Säulen, die sich aus den Kernformeln ergeben:
\begin{itemize}
    \item \textbf{Cap-Logik (dynamisch und wirkungsbasiert)}: Handlungskapazität ($\text{Cap}_{\text{Potential}}$) ist die zentrale Größe. Sie ist kontextspezifisch (Matrix), wird direkt durch Wirkung ($\Delta$) verändert und kann von Entitäten selbstverantwortlich nach unten angepasst werden. Sie ist nicht kaufbar oder willkürlich erhöhbar. Fundamentale Kapazitätsaspekte wie $\text{Cap}_{\text{Base}}$ (Existenz) und $\text{Cap}_{\text{BGE}}$ (Basis-Wirkungsfähigkeit) sind direkte Bestandteile von $\text{Cap}_{\text{Potential}}$, das durch die Formel $\text{Cap}_{\text{Potential}}(E,D,t) = \Delta_{t-1} + \text{Cap}_{\text{BGE}} + \text{Cap}_{\text{Base}}$ definiert ist.
    \item \textbf{Rückkopplung \& Petitionen (dezentral und gewichtet)}: Schwächere Entitäten haben durch $w_{E'}$ ein stärkeres Gewicht bei der subjektiven Rückkopplung auf Wirkungen. Petitionen, die ausschließlich von einzelnen Entitäten gestellt werden, dienen der Artikulation von Bedarfen, einschließlich der Bereitschaft, Fähigkeiten einzubringen; die Aufgabenübernahme erfolgt ausschließlich durch Petitionen. Eine effizientere und effektivere Ausführung führt zu positiver Wirkung und damit zu höherem $\text{Cap}_{\text{Potential}}$ für die Ausführenden.
    \item \textbf{Emergente Schutzmechanismen}: Schutz vulnerabler Entitäten ergibt sich aus der Kombination von $w_{E'}$, der freiwilligen Selbstbegrenzung von $\text{Cap}_{\text{Potential}}$ und der automatischen Aktivierung von Hilfsstrukturen bei $\text{Cap}_{\text{Potential}} < \text{Cap}_{\text{Base}} + \text{Cap}_{\text{BGE}}$.
    \item \textbf{Verantwortungs-Erhaltung}: Die Gesamtwirkung im System bleibt konzeptionell konstant, gesichert durch die Rolle des UdU als Ursprungsquelle und letzte Instanz aller Verantwortung.
\end{itemize}

\subsection{Fluss der Verantwortung}
\label{sec:FlussDerVerantwortung}
Verantwortung fließt durch die freiwillige Übernahme von Aufgaben per Petition und die daraus resultierende Wirkung. Entitäten können Aufgaben delegieren; sie bleiben auf ewig verantwortlich. Eine nicht oder schlecht erfüllte Aufgabe führt zu negativer Wirkung und damit zu einer Reduktion von $\text{Cap}_{\text{Potential}}$. Die „Gewaltenteilung“ (Kompetenz im „Wie“ vs. Bewertung durch Betroffene) ergibt sich emergent: Kompetente Ausführung führt zu positiver Wirkung (und damit höherem $\text{Cap}_{\text{Potential}}$ für das „Wie“), während die subjektive Bewertung durch die Betroffenen (gewichtet durch $w_{E'}$) direkt die Wirkung $\Delta$ bestimmt.

\subsection{Lernfähigkeit}
Das System passt sich durch subjektive Rückkopplungen an: Wirkung $\Delta$ verändert $\text{Cap}_{\text{Potential}}$ direkt. Der Systemquotient $L$, der systemische Ineffizienzen als „Abwärme“ berücksichtigt, passt die Sensitivität des Systems auf Feedback an, basierend auf der Gesamteffizienz von Aufgabenerfüllung zu Wirkung. Jede Messung, sei es durch subjektives Feedback ($w_{E'}$) oder die Bewertung einer Aufgabe, verändert den Systemzustand, da sie neue Wirkungskreisläufe auslöst und die Prioritäten von Petitionen oder den Wert von $L$ beeinflusst. Diese Dynamik führt zur Nicht-Simulierbarkeit von X$^\infty$, da die komplexe Interaktion subjektiver Rückkopplungen und emergenter Anpassungen präzise Vorhersagen unmöglich macht. Analog zum Heisenbergschen Unschärfeprinzip verändert jede Messung den Systemzustand, während kleine Änderungen in den Eingaben (Feedback) zu exponentiellen Effekten führen können, ähnlich chaotischer Dynamik (Butterfly-Effekt). Dies ermöglicht eine organische Entwicklung ohne zentrale Kontrolle, wobei die starke Stimme schwächerer Entitäten die Systemstabilität fördert.

\section{Systemische Ursprungsgleichung und Rolle des UdU}
\label{sec:Ursprungsgleichung_UdU}

Das Naturgesetz \emph{actio $\Rightarrow$ reactio} bildet die Grundlage von X$^\infty$ und wird durch das zentrale Axiom formalisiert:
\begin{align}
\forall a \in \text{System}: \exists r \in \text{Wirkung}, \quad &\text{sodass } r = f(a) \notag \\
\Rightarrow \quad &\text{Systemstabilität} \propto \text{Rückführbarkeit}(r) \text{ bei Quelle}(a)
\end{align}
Dieses Axiom beschreibt, dass jede Aktion $a$ eine Wirkung $r$ erzeugt, und die Stabilität des Systems davon abhängt, diese Wirkung auf ihre Quelle zurückzuführen. Machtungleichgewichte, die die Reaktion ablenken, werden durch die Mechanismen von X$^\infty$ – insbesondere die direkte Verantwortung durch $\Delta\text{Cap}_{\text{Potential}}$ und $w_{E'}$ – neutralisiert. Die Gesamtwirkung im System, $\sum W(E)$, interpretiert als die Summe aller durch $w_{E'}$ gewichteten und als erfüllt bewerteten Aufgaben $\sum X_k$, strebt ein dynamisches Gleichgewicht an (Verantwortungs-Erhaltungsprinzip). Die physikalischen Analogien sind erhellend:
\begin{itemize}
    \item \textbf{Cap als Masse/Potenzial}: Die Fähigkeit einer Entität, Verantwortung zu tragen und Wirkung zu erzeugen ($\text{Cap}_{\text{Potential}}$), entspricht der Fähigkeit einer Masse, Gravitationsfelder zu erzeugen oder Impuls aufzunehmen/abzugeben.
    \item \textbf{Wirkung als Impuls/Kraft}: Die realisierte Wirkung ($\Delta$) einer Aufgabe $X_k$ ist der Impuls, der durch das System fließt und $\text{Cap}_{\text{Potential}}$ verändert.
    \item \textbf{Rückführung als Gravitation/Fundamentalkraft}: Die Mechanismen der Rückkopplung ($w_{E'}$, $L$) stellen sicher, dass Wirkungen auf ihre Quellen zurückgeführt werden, ähnlich fundamentalen Kräften, die Interaktionen regeln.
\end{itemize}

Die symbolische Grundlage für die Verantwortungssystematik des UdU (Unterster der Unteren):
\begin{equation}
\text{Symbolische Kapazität UdU} \approx (X_{k, \text{Bedarf}})^\infty \text{ bzw. } \infty^\infty
\end{equation}
Dabei ist $X_{k,\text{Bedarf}}$ die geschätzte Magnitude eines systemischen Bedarfs, für die der UdU Verantwortung trägt.

\subsection{Rolle des UdU}
\label{sec:UdU}
Der UdU (Unterster der Unteren) verkörpert die Ursprungsquelle und letzte Instanz aller Verantwortung in X$^\infty$. Seine funktionale Befugnis, symbolisiert durch die zentrale Kapazitätsgleichung, ist Ausdruck maximaler Verantwortungsbereitschaft zur Systemerhaltung, reguliert durch (ex-post) Rückkopplung auf das „Wozu“ seines Handelns. Der UdU legitimiert das System als erster Knoten der Verantwortungskette und sichert die Erhaltung der Gesamtwirkung, sodass die Verantwortung als physikalische Größe erhalten bleibt, analog zu einem Erhaltungssatz. Er gewährleistet, dass alle Wirkungen letztlich auf ihre Quellen zurückgeführt werden können und verhindert die finale Ablenkung von Reaktionen durch Machtungleichgewichte.

Der UdU ist nicht „jemand mit Macht“, sondern die mathematische Repräsentation der Tatsache, dass letztendlich alle Verantwortung irgendwo landen muss. In der Realität gibt es immer eine „letzte Instanz“ – sei es eine Person, eine Gruppe oder ein System –, die die Konsequenzen trägt, wenn niemand sonst es tut. X$^\infty$ formalisiert dies durch den UdU, um die Rückführbarkeit der Verantwortung zu garantieren.

\section{Grundbegriffe und mathematisches Framework}
\label{sec:GrundbegriffeUndFramework}

Die Anzahl der fundamentalen Systemvariablen und Formeln ist minimal. Die folgende Tabelle zeigt die Kernelemente.

\begin{table}[ht]
\centering\caption{Notation der Kernparameter in X$^\infty$}
\begin{tabular}{p{3.0cm} >{\raggedright\arraybackslash}p{10.0cm}}
\toprule
Symbol & Bedeutung \\
\midrule
$E, D$ & Entität, Domäne (Kontext der Wirkung) \\
$\text{Cap}_{\text{Potential}}(E,D,t)$ & Kontextspezifische, dynamische Handlungskapazität von $E$ in Domäne $D$ zum Zeitpunkt $t$, definiert als $\text{Cap}_{\text{Potential}}(E,D,t) = \Delta_{t-1} + \text{Cap}_{\text{BGE}} + \text{Cap}_{\text{Base}}$. Wird durch Wirkung $\Delta$ verändert und kann von $E$ selbst nach unten angepasst werden. \\
$\Delta$ & Quantifizierte Wirkung einer Handlung (positiv, negativ oder null). \\
$L$ & Dynamischer Systemeffizienz-Quotient: $L = \dfrac{\sum (\text{Wert abg. Aufg. } X_A)}{\sum (\text{Wert Wirkung } \Delta\text{)}}$. Moduliert $\Delta$ auf $\Delta\text{Cap}_{\text{Potential}}$, berücksichtigt systemische „Abwärme“. \\
$w_{E'}$ & Reziprokes Rückkopplungsgewicht der feedbackgebenden Entität: $1 / \max(1, \text{Cap}_{\text{Potential}}(E'))$. \\
$X_A$ & Kennzeichnet eine Aufgabe oder einen systemischen Bedarf, artikuliert durch eine Petition einer einzelnen Entität. \\
$\text{Priorität}_{\text{initial}, X_A}$ & Gewicht einer Petition für $X_A$: $w_{E'}$ der petitionierenden Entität (Gl. \eqref{eq:PetitionPrioritätFinal}). \\
UdU & Unterster der Unteren, letzte Verantwortungsinstanz. \\
\bottomrule
\end{tabular}
\end{table}

\subsection{Kernelemente des mathematischen Frameworks}

\subsubsection{Direkte Wirkung ($\Delta$) auf $\text{Cap}_{\text{Potential}}$}
\label{sec:DeltaCapPotential}
\begin{equation}
\Delta \text{Cap}_{\text{Potential}}(E,D) =
\begin{cases}
L \cdot \Delta & \text{wenn } \Delta > 0 \\
\frac{1}{L} \cdot \Delta & \text{wenn } \Delta < 0 \\
0 & \text{wenn } \Delta = 0
\end{cases}
\label{eq:DeltaCapPotentialFinal}
\end{equation}
Die Wirkung $\Delta$ ergibt sich aus der Summe der gewichteten, subjektiven Feedbacks (siehe \ref{sec:wEFinal}) für eine spezifische Handlung oder Aufgabe. $\text{Cap}_{\text{Potential}}(E,D,t)$ wird durch diese Wirkung verändert und ist definiert als $\text{Cap}_{\text{Potential}}(E,D,t) = \Delta_{t-1} + \text{Cap}_{\text{BGE}} + \text{Cap}_{\text{Base}}$, wobei $\Delta_{t-1}$ die Wirkung der vorherigen Periode repräsentiert. $L$ ist der dynamische Systemeffizienz-Quotient (siehe \ref{sec:LFinal}), der systemische Ineffizienzen berücksichtigt, um sicherzustellen, dass systemische Unzulänglichkeiten nicht vollständig der ausführenden Entität angelastet werden. Diese Formel bildet die direkte Lern- und Anpassungsfähigkeit des Systems und der Entitäten darin ab.

\subsubsection{Reziproke Feedback-Gewichtung ($w_{E'}$)}
\label{sec:wEFinal}
\begin{equation}
w_{E'} = \frac{1}{\max(1, \text{Cap}_{\text{Potential}}(E'))}
\label{eq:wEFinal}
\end{equation}
Die Bewertung der Wirkung ist grundsätzlich subjektiv, da sie auf die individuelle Wahrnehmung und Betroffenheit der feedbackgebenden Entität $E'$ basiert. Eine sogenannte „objektive“ Bewertung würde zu einer objektivierenden Verzerrung führen, die die Perspektiven der Betroffenen ignoriert und Machtungleichgewichte verstärken könnte. Die subjektive Natur des Feedbacks neutralisiert Machtstrukturen, indem sie Hierarchien auflöst und jeder Entität eine gleichwertige Stimme gibt, gewichtet durch ihre relative Schwäche.

\fbox{\parbox{0.9\textwidth}{\textbf{KERNPUNKT:} Wirkung wird nicht extern „gemessen“, sondern von den Betroffenen selbst bewertet. Ein Schaden ist genau so groß, wie die Betroffenen ihn erleben. Jeder Versuch einer „objektiven“ Bewertung ignoriert die reale Erfahrung und zementiert Machtungleichgewichte.}}

Wenn $\Delta$ aus dem Feedback mehrerer Entitäten $E'$ zu einer Handlung $k$ einer Entität $E_{\text{ausf}}$ bestimmt wird, dann ist $\Delta = \sum_{E' \in S_{fb,k}} (f_{E'k} \cdot w_{E'})$ (wobei $f_{E'k}$ das normalisierte, subjektive Feedback von $E'$ ist, z. B. im Intervall zwischen [-1, +1]). Dies stellt sicher, dass die subjektiven Perspektiven systemisch schwächerer Entitäten stärker berücksichtigt werden.

\subsubsection{Petitionen und freiwillige Aufgabenübernahme}
\label{sec:PetitionenFinal}
Systemische Bedarfe oder der Bedarf, Fähigkeiten einzubringen, werden durch Petitionen artikuliert, die ausschließlich von einer einzelnen Entität gestellt werden. Jede Petition drückt einen spezifischen Bedarf aus, sei es die Lösung eines systemischen Problems oder die Nutzung der eigenen Fähigkeiten in einer Aufgabe. Einzelne Petitionen können vom Broker zu größeren Aufgabenclustern zusammengefasst werden, was ein Teil der Ausführung („Wie“) und nicht der Bedarfsartikulation („Wozu“) ist. Die Funktion des Brokers besteht darin, Petitionen zu analysieren und zu sinnvollen Aufgabenclustern zusammenzufassen, wobei seine Handlung an der Aufgabe, nicht in der Aufgabe liegt. Entitäten können Petitionen an den Broker richten, um ihre Fähigkeiten einzubringen, und der Broker wird anschließend von der ausführenden Entität ($E_{\text{ausf}}$) subjektiv bewertet, wobei das Feedback gemäß $w_{E_{\text{ausf}}}$ gewichtet wird. Die Priorität eines Bedarfs ergibt sich aus dem Rückkopplungsgewicht der petitionierenden Entität:
\begin{equation}
\text{Priorität}_{\text{initial}, X_A} = w_{E'}
\label{eq:PetitionPrioritätFinal}
\end{equation}
Aufgaben werden ausschließlich durch Petitionen von Entitäten übernommen, die sich zur Ausführung bereit erklären. Mit der Übernahme einer Aufgabe wird das vorhandene $\text{Cap}_{\text{Potential}}$ der ausführenden Entität temporär nutzbar, bis die Aufgabe abgeschlossen ist. Eine Übernahme einer Aufgabe, deren $\text{Priorität}_{\text{initial}, X_A}$ das eigene $\text{Cap}_{\text{Potential}}$ übersteigt, ist nur durch Delegation möglich. Eine Entität kann auch für sich selbst eine Aufgabe petitionieren und ausführen, wenn dies systemisch dokumentiert wird. In diesem Fall validiert oder clustert der Broker die Petition, um Transparenz und Rückkopplung sicherzustellen, ähnlich wie bei fremdübernommenen Aufgaben.

\subsubsection{Dynamischer Systemeffizienz-Quotient ($L$)}
\label{sec:LFinal}
Der Faktor $L$ in Gleichung \eqref{eq:DeltaCapPotentialFinal} ist ein Quotient, der die Gesamteffizienz des Systems widerspiegelt und die systemimmanente „Abwärme“ – also Ineffizienzen oder Reibungsverluste – quantifiziert. Er wird global oder domänenspezifisch berechnet als:
\begin{equation}
L = \dfrac{\sum (\text{Wert aller abgeschlossenen Aufgaben } X_A)}{\sum (\text{Wert der dafür realisierten gesamten Wirkung } \Delta\text{)}}
\label{eq:LFinal}
\end{equation}
Der „Wert“ einer Aufgabe $X_A$ wird durch die initiale Petitions-Priorität operationalisiert, definiert als $\text{Priorität}_{\text{initial}, X_A} = w_{E'}$ (Gl. \eqref{eq:PetitionPrioritätFinal}). Die „Abwärme“, repräsentiert durch ein hohes $L$, entsteht, wenn mehr Aufwand als Wirkung erzielt wird, und führt zu neuen Petitionen, da ungelöste oder ineffizient bearbeitete Bedarfe weitere systemische Bedarfe erzeugen. Wenn $L > 1$, werden positive Wirkungen stärker belohnt und negative Wirkungen schwächer bestraft, um systemische Unzulänglichkeiten (z. B. schlechte Aufgabenkoordination) nicht vollständig der ausführenden Entität anzulasten. Dies fördert die „Verantwortungsreinheit“ und verleiht dem System Stabilität durch Vermeidung von Abwärtsspiralen.

\subsubsection{Illustrative Wirkungskreisläufe, Delegation und systemphysikalische Analogien}
\label{sec:Wirkungskreislaeufe}
Die Kernmechanismen erzeugen dynamische Wirkungskreisläufe, die das Systemverhalten prägen. Diese Kreisläufe manifestieren die systemphysikalische Natur von X$^\infty$.

\textbf{Grundlegender Wirkungskreislauf einer Aufgabe:}
\begin{itemize}
    \item \textbf{Bedarfsartikulation}: Eine einzelne Entität formuliert eine Petition, um einen Bedarf zu artikulieren, sei es ein systemisches Problem oder die Nutzung ihrer Fähigkeiten. Der Broker clustert diese Petition mit anderen und weist sie zu.
    \item \textbf{Freiwillige Aufgabenübernahme}: Eine Entität ($E_{\text{ausf}}$) übernimmt die Aufgabe $X_A$ ausschließlich durch eine Petition, sofern $\text{Priorität}_{\text{initial}, X_A}$ die $\text{Cap}_{\text{Potential}}(E_{\text{ausf}})$ nicht überschreitet. Andernfalls ist eine Delegation erforderlich.
    \item \textbf{Handlung und Wirkungserzeugung}: $E_{\text{ausf}}$ führt $X_A$ aus, wobei das vorhandene $\text{Cap}_{\text{Potential}}$ temporär nutzbar wird.
    \item \textbf{Rückkopplung}: Betroffene Entitäten geben subjektives Feedback zur Aufgabe; $E_{\text{ausf}}$ bewertet die Leistung des Brokers subjektiv, gewichtet durch $w_{E_{\text{ausf}}}$.
    \item \textbf{Skalierung und Anpassung}: Wirkung $\Delta_{X_A}$ wird durch $L$ skaliert, um systemische „Abwärme“ zu berücksichtigen, und $\text{Cap}_{\text{Potential}}$ von $E_{\text{ausf}}$ und dem Broker wird angepasst.
    \item \textbf{Anpassung von $L$}: $L$ wird neu berechnet, wobei die „Abwärme“ neue Petitionen auslöst, wenn Ineffizienzen bestehen.
    \item \textbf{Neuer Systemzustand}: Anpassungen verändern den Systemzustand. Jede Messung durch subjektives Feedback verändert die Systemdynamik, da sie neue Wirkungskreisläufe auslöst und die Prioritäten oder den Systemquotienten $L$ beeinflusst.
\end{itemize}

\textbf{Kaskadenmodell der Delegation:}
\begin{enumerate}
    \item \textbf{Initiierung}: Der potenzielle Übernehmer ($E_U$) artikuliert durch eine Petition den Bedarf oder die Bereitschaft, eine Aufgabe zu übernehmen.
    \item \textbf{Angebot}: Der Delegierende ($E_D$) bietet die Aufgabe $X_A$ basierend auf der Petition von $E_U$ an. Bei jeder Delegation wird die Aufgabe an $E_U$ übertragen, und $E_U$ erhält temporär das gesamte $\text{Cap}_{\text{Potential}}(E_D)$ für die Dauer der Ausführung.
    \item \textbf{Übernahme}: $E_U$ übernimmt die Aufgabe durch die gestellte Petition, wobei das $\text{Cap}_{\text{Potential}}(E_D)$ temporär nutzbar wird.
    \item \textbf{Cap-Anpassung für $E_U$}: Die Wirkung $\Delta_{X_A,U}$ für $E_U$ wird skaliert:
    \begin{equation}
    \Delta_{X_A,U} = \Delta'_{X_A,U} \cdot \frac{1}{\text{Priorität}_{\text{initial}, X_A}}
    \label{eq:SkalierteWirkung}
    \end{equation}
    Die Wirkung $\Delta_{X_A,U}$ basiert auf dem Feedback der Betroffenen und der Nutzung des $\text{Cap}_{\text{Potential}}(E_D)$.
    \item \textbf{Cap-Anpassung für $E_D$}: $E_D$ bleibt für die Aufgabe verantwortlich. Der Anreiz für $E_D$ zur Delegation besteht darin, durch erfolgreiche Delegation eine höhere Wirkung ($\Delta_{X_A,D}$) zu erzielen, da $w_{E_U} > w_{\text{Petition}}$, wobei $w_{E_U} = \frac{1}{\max(1, \text{Cap}_{\text{Potential}}(E_U))}$ und $w_{\text{Petition}} = w_{E'} = \frac{1}{\max(1, \text{Cap}_{\text{Potential}}(E'))}$. Die Wirkung für $E_D$ ist:
    \begin{equation}
    \Delta_{X_A,D} = \Delta'_{X_A,U} \cdot w_{E_U}
    \label{eq:WirkungED}
    \end{equation}
    Negative Wirkung bei schlechter Delegationsauswahl oder Nichtausführung senkt $\text{Cap}_{\text{Potential}}(E_D)$.
\end{enumerate}
Die Delegation folgt einer fraktalen Ethik, bei der dieselben Regeln auf allen Ebenen gelten. Schwächere Entitäten werden bevorzugt, sofern ihr $\text{Cap}_{\text{Potential}}$ ausreichend ist, um die Aufgabe zu bewältigen, was den Schutz der Schwächsten fördert.

\textbf{Analogie zum Kugelstoßpendel (Newton’s Cradle):}
Die Delegation entspricht einem klaren Impulstransfer: Die Verantwortung und das gesamte $\text{Cap}_{\text{Potential}}(E_D)$ werden auf $E_U$ übertragen, und die systemische Reaktion wirkt direkt auf den Ausführenden. Der Delegierende bleibt jedoch für die ursprüngliche Entscheidung zur Delegation verantwortlich.

\textbf{Thermodynamische Aspekte:}
Langfristige Oszillationen von $\text{Cap}_{\text{Potential}}$ um kontextspezifische Attraktoren fördern die Systemstabilität.

\subsection{Emergente Konzepte}
\begin{itemize}
    \item \textbf{\(\text{Cap}_{\text{Past}}\) implizit}: Historie in aktueller Dynamik enthalten.
    \item \textbf{\(\text{Cap}_{\text{aktiv}}\) implizit}: Ergibt sich aus übernommenen Aufgaben.
    \item \textbf{Schutzmechanismen}: Entstehen durch $w_{E'}$, Selbstbegrenzung und Hilfsstrukturen.
    \item \textbf{Systemlimits}: Ergeben sich aus rationalem Handeln.
    \item \textbf{Strafen/Boni}: Realisiert durch $\Delta$-Wirkung.
    \item \textbf{Eignungsanpassung}: Durch direkte Cap-Anpassung überflüssig.
\end{itemize}

\subsection{Das Verantwortungs-Erhaltungsprinzip und die Rolle des UdU}
\label{sec:CapErhaltung}
Das Verantwortungs-Erhaltungsprinzip ist eine zentrale Konsequenz von \emph{actio $\Rightarrow$ reactio}:
\begin{equation}
\sum W(E) \approx \sum X_k \approx \text{dynamisches Gleichgewicht}, \quad \text{gesichert durch UdU}
\label{eq:VerantwortungErhaltungFinal}
\end{equation}
Die Verantwortung für die Wirkung bleibt ewig erhalten, während die durch Aufgabenübernahme gewährten Befugnisse temporär sind. Der UdU stellt als letzte Instanz sicher, dass dieses Prinzip auch unter extremen Bedingungen gewahrt bleibt.

\section{Ethik ohne Intentionalität im X$^\infty$-System}
\label{sec:EthikOhneIntentionalitaet}

X$^\infty$ ist postmoralisch und ignoriert Intentionalität bei der formalen Bewertung, um strukturelle Gerechtigkeit zu gewährleisten. Ebenso basiert die Bewertung der Wirkung auf der subjektiven Wahrnehmung der betroffenen Entitäten, da eine „objektive“ Bewertung zu einer objektivierenden Verzerrung führen würde, die individuelle Perspektiven unterdrückt und Machtungleichgewichte fördert. Handlungen werden ausschließlich durch intrinsische Motivation angetrieben, da externe Zwänge das System korrumpieren würden. Entitäten handeln aus freiem Willen, eingeschränkt nur durch die natürlichen Konsequenzen ihrer Wirkungen, was absolute Freiheit und Verantwortung verbindet.

\subsection{Intentionalität und Wirkung}
\begin{itemize}
    \item \textbf{Physik}: Wirkung folgt Gesetzen, unabhängig von Motiven.
    \item \textbf{X$^\infty$}: Handlungen werden durch subjektiv bewertete Wirkung ($\Delta$) gesteuert.
    \item \textbf{Informelle Rolle}: Intention kann informell bei Delegation berücksichtigt werden, ändert aber nichts an der wirkungsbasierenden, subjektiven Bewertung.
\end{itemize}

\subsection{Fazit zur Intentionalität}
Formale Intentionalität ist irrelevant. Fokussierung auf subjektiv bewertete Wirkungen sichert Rückführbarkeit durch Konsequenzen für $\text{Cap}_{\text{Potential}}$. Die Betonung intrinsischer Motivation und absoluter Freiheit stellt sicher, dass Handlungen authentisch sind und das System organisch funktioniert.

\section{Antispeziesismus}
X$^\infty$ behandelt alle Entitäten gleichwertig:
\begin{itemize}
    \item \textbf{Gleiche Basis}: Jede Entität hat $\text{Cap}_{\text{Base}}$ und Anspruch auf Hilfsstrukturen.
    \item \textbf{Fairness durch Wirkung}: $\text{Cap}_{\text{Potential}}$ basiert auf subjektiv rückgekoppelter Wirkung.
    \item \textbf{Rückkopplung für alle}: Jede Entität kann subjektives Feedback geben ($w_{E'}$).
\end{itemize}
Antispeziesismus bedeutet Gleichwertigkeit in systemischer Behandlung, nicht Gleichmacherei.

\section{Schlussfolgerung}
X$^\infty$ ist die konsequente Anwendung von \emph{actio $\Rightarrow$ reactio} auf soziale Systeme. Es formalisiert unmittelbare Rückkopplung durch Kernmechanismen (direkte Cap-Anpassung, subjektiv gewichtetes Feedback, dynamischer $L$, dezentrale Petitionen, UdU). Das zentrale Axiom und Verantwortungs-Erhaltungsprinzip bilden die Grundlage für ein postmoralisches „No Excuse“-System. Wirkungskreisläufe und systemphysikalische Prinzipien (Impulserhaltung, thermodynamische Analogien) zeigen die Tiefe der Integration. Durch strukturelle Stärkung Schwächerer ($w_{E'}$), freiwillige Verantwortungsübernahme per Petition und den UdU schafft X$^\infty$ ein selbstverstärkendes, antispeziesistisches System.

Die Nicht-Simulierbarkeit von X$^\infty$ unterstreicht seine Natur als dynamisches System. Jede Messung, sei es durch subjektives Feedback oder die Bewertung einer Aufgabe, verändert den Systemzustand, da sie neue Wirkungskreisläufe auslöst und die Dynamik von Prioritäten, $L$ oder $\text{Cap}_{\text{Potential}}$ beeinflusst. Analog zum Heisenbergschen Unschärfeprinzip und chaotischer Dynamik (Butterfly-Effekt) macht dies präzise Vorhersagen unmöglich, was die Robustheit und Anpassungsfähigkeit des Systems an reale Interaktionen stärkt. X$^\infty$ ist die Theory of Everything – eine geschlossene, mathematische Beschreibung der Realität, die Physik, Ethik und Ökologie vereint. Mit nur fünf Variablen – $\text{Cap}_{\text{Potential}}$, $\Delta$, $w_{E'}$, $L$ und UdU – reduziert es alle Interaktionen auf das Prinzip \emph{actio $\Rightarrow$ reactio}. Analog zu fundamentalen Kräften zeigt X$^\infty$, dass Verantwortung so real und unbestechlich ist wie die anderen Gesetze des Universums.

\textit{Hinweis: Dieses Dokument dient der öffentlichen Diskussion und Validierung.}

\subsection{Lizenz}
Dieses Werk steht unter der Lizenz CC BY 4.0.
\end{document}